\subsection{Further Aspects of Key Management}\label{sec:further:key_management}

The security of cryptographic keys is not limited to just a proper definition of the functional aspects of key management (i.e., those described in \Cref{sec:key_management}) and the choice of the most suitable Cloud-based \gls{kms} services (see \Cref{sec:cloud}). Therefore, below we identify and briefly discuss further aspects relevant to the secure development and deployment of key management and \glspl{kms}.



\subsubsection{Policies}\label{sec:further:key_management:policies}

The management of cryptographic keys involves the definition, enforcement, and continuous revision of various policies --- a policy can be seen as a set of rules pertaining to a specific aspect of a system (e.g., a \gls{kms}). Organizations usually structure policies in a hierarchy, where each policy in a certain layer details the enforcement and implementation of the rules specified in the policies of the higher layer(s) \cite{barker_framework_2013}. In other words, top-level policies define security, business, and regulatory requirements deriving from the organizations themselves, best practices, and relevant standards (e.g., \gls{nist}'s), while intermediate- and low-level policies describe --- in greater detail --- suitable strategies and procedures to achieve such requirements. For instance, a high-level policy requiring to ``protect sensitive data'' may be linked to an intermediate-level policy requiring to ``encrypt sensitive data in transit and at rest'' which, in turn, may be linked to a low-level policy requiring to ``encrypt data in transit using \gls{tls} and data at rest using \gls{aes}-128 in a validated \gls{fips} 140-3 Hardware-based cryptographic module''. Policies at higher levels should be easily understandable by users (e.g., textual documents, tables, flow charts), while policies at lower levels may also be directly enforceable by automated systems. In any case, policies should always be written in a formal language for unambiguous meaning. In the context of key management, \gls{nist} defines a hierarchy of policies consisting of 3 layers \cite{barker_framework_2013}:

\begin{itemize}
	\item \textbf{information management policy}: this high-level policy expresses the security, business, and regulatory requirements of the organization on the use of (various categories of) sensitive data, the definition of responsibilities and management roles, and the need for security properties such as confidentiality, integrity, and availability. The information management policy also considers industrial best practices, relevant standards, and regulatory frameworks, and it is often complemented with a risk assessment \cite{joint_guide_2012};
	\item \textbf{information security policy}: this intermediate-level policy refines the aforementioned requirements by specifying how data is categorized, the sensitivity levels assigned to data, eventual access constraints, and the permissions to distribute to the management roles. The information security policy also identifies protections for data (e.g., encryption), taking into account the threats identified in the risk assessment previously conducted;
	\item \textbf{\gls{ckmssp}}: this low-level policy --- previously known as ``Key Management Policy'' --- specifies rules for the management of keys and metadata in each phase of the lifecycle shown in \Cref{fig:lifecycle} (e.g., installation/registration, generation, utilization, rotation), with the final goal to preserve the confidentiality, integrity, availability, and authenticity of keys and metadata \cite{barker_recommendation_2019_p2}. Intuitively, the rules within the \gls{ckmssp} need to be consistent with higher-level policies and have a clear mapping with the requirements defined by the organization. Moreover, the chosen \gls{kms} --- regardless of whether it was developed by the organization itself or an external party --- should be able to enforce the \gls{ckmssp}. Usually, a \gls{ckmssp} is composed of a policy for each phase of the key management lifecycle (e.g., key generation policy, key usage policy, key recovery policy). Furthermore, a \gls{ckmssp} is usually complemented with (at least) the following policies:
	\begin{itemize}
		\item \textbf{certificate policy}: this policy regulates all aspects associated with the management (e.g., issuance, distribution, revocation, renewal) of digital certificates, including details related to their types, lifecycle, and domains of applicability. Furthermore, a certificate policy defines the architecture and the entities that compose the \gls{pki} (e.g., accredited \glspl{ca}), their roles, and their duties. Finally, a certificate policy identifies which cryptographic ciphers to use (and their configuration) for digital certificates. For more details on certificate policies, we refer the interested reader to \cite{ford_internet_2003};
		\item \textbf{key-inventory policy}: this policy outlines guidelines and procedures for managing and maintaining an inventory of cryptographic keys (see \Cref{sec:key_management:inventory}). A key-inventory policy defines what keys and metadata should be stored in the inventory management system \cite{barker_recommendation_2020_p1}, including information on responsibilities, monitoring (e.g., for expired keys and insecure cryptographic parameters), and any specific requirements imposed by legal or regulatory frameworks. For more details on key-inventory policies, we refer the interested reader to \cite{barker_recommendation_2019_p2};
		\item \textbf{\gls{ac} policy}: this policy --- usually part of \gls{iam} (see S8-CSC-D1 in \Cref{sec:intro:scope}) --- comprises rules stating which agents can perform which actions (e.g., call for key rotation, encryption) on which of the resources (e.g., keys) maintained by the \gls{kms}. A formal representation of these rules is provided by the \gls{ac} model --- e.g., \gls{rbac} \cite{sandhu_access_1998} and \gls{abac} \cite{horwitz_toward_2002} are two models --- giving the semantics for granting or denying agents’ requests. The hardware and software entities that enforce the \gls{ac} policy based on the chosen model compose the so-called ``enforcement''. The enforcement may be either centralized --- e.g., based on a standard like \gls{xacml}\footnote{OASIS eXtensible Access Control Markup Language (XACML) TC (\url{https://www.oasis-open.org/committees/tc_home.php?wg_abbrev=xacml})} or monitors like the \gls{opa}\footnote{Open Policy Agent (\url{https://www.openpolicyagent.org/})} --- or decentralized --- hence, enforced directly through cryptographic with a \gls{cac} scheme like those in \cite{berlato_formal_2022,berlato_end_2022}. For more details on \gls{ac} policies, we refer the interested reader to \cite{nist_security_2020};
		\item \textbf{cryptographic algorithm policy}: this policy selects cryptographic algorithms, ciphers, and cryptosystems, as well as key lengths approved for use in the \gls{kms}. This selection is usually based on the requirements of the organization, industrial standards, and best practices. For more details on the cryptographic algorithm policy, see \Cref{app:ciphers};
	\end{itemize}
\end{itemize}

As a final remark, organizations operating in different scenarios (e.g., financial services, pharmaceuticals, government bodies) are likely to have different policies and rules. Furthermore, cross-organizational scenarios may expect the definition of cross-domain security policies for interoperability \cite{barker_framework_2013}.



\subsubsection{Authentication and Authorization}\label{sec:further:key_management:authn}

In line with the Zero Trust security model\footnote{Zero Trust - Glossary | CSRC (\url{https://csrc.nist.gov/glossary/term/zero\_trust})} (``never trust, always verify''), any entity performing any operation within the key management lifecycle should be properly authenticated and authorized \cite{chandramouli_cryptographic_2013}. Authentication is the process of verifying --- after having previously registered --- the (digital) identity of users, processes, devices, and agents in general when attempting to access a resource or perform a digital transaction (for more details on authentication, see S8-CSC-D1 in \Cref{sec:intro:scope}), while authorization is the process of verifying to what resources (and actions) a (usually authenticated) agent has access \cite{samarati_access_2001}; the authorization process often consists of evaluating \gls{ac} rules contained in a policy to make and then enforce \gls{ac} decisions --- e.g., grant or deny an access request. In the context of cryptographic key management, authentication and authorization are of fundamental importance. For instance, the association between a digital identity and its owner is crucial for maintaining a record of activities and data access history (see \Cref{sec:further:key_management:logging}) to guarantee accountability and identify malicious activities. Moreover, the assignment of permissions to agents --- where a permission represents the ability to perform an action over a resource (e.g., rotate a cryptographic key) --- allows for regulating (and preventing unauthorized) access, checking that all applicable constraints are satisfied, and enforcing the least privilege principle;\footnote{least privilege - Glossary | CSRC (\url{https://csrc.nist.gov/glossary/term/least\_privilege})} \gls{cac} is the natural combination of cryptographic and \gls{ac} \cite{berlato_formal_2022,berlato_end_2022}.

\sbnote{TODO talk also about SoD
%Separation of Duties: Separation of duties 
%By creating distinct duties or user types, that can only perform certain operations, there is an additional layer of security applied, since multiple users would need to be in collusion for key material to be stolen or manipulated in an inappropriate manner, and the audit trail to be erased
}

\subsubsection{Logging and Auditing}\label{sec:further:key_management:logging}

In a \gls{kms} --- and key management in general --- any event concerning or relevant to cryptographic keys (i.e., creation, access, use, transmission, creation of backups, transition among key states and phases) should be properly recorded in a log and, when necessary, in the metadata of keys as well \cite{barker_profile_2015}. Logs are usually collected in a distributed fashion but then stored in a centralized way; this approach implies both the ability to integrate logs coming from different sources (e.g., the Cloud, multiple \glspl{kms}) and also filter and review logs --- which may be voluminous --- with automated assistance \cite{csa_key_2020}. Of course, the ability to preserve and demonstrate the integrity of logs is of the utmost importance for the following use; indeed, an organization can use logs for several purposes:

\begin{itemize}
	\item \textbf{security monitoring}: analyze events and operations (preferably in real-time) to detect anomalies (e.g., unusual patterns or suspicious activities) and attacks (e.g., unauthorized access, multiple failed login attempts);
	\item \textbf{incident response}: review logs to identify the root cause of security incidents, conduct forensic analyses, reconstruct the sequence of actions leading to security incidents, and develop new effective security measures;
	\item \textbf{compliance}: demonstrate compliance with regulations --- e.g., \gls{gdpr},\footnote{GDPR compliance (\url{https://gdpr.eu/})} \gls{hipaa},\footnote{HIPAA Home | HHS.gov (\url{https://www.hhs.gov/hipaa/index.html})} and \gls{picdss}\footnote{PCI Security Standards Council (\url{https://www.pcisecuritystandards.org/})} --- by allowing third-party auditors to review the activities of a system and identify behaviors not compliant with the established operational policies and procedures (see \Cref{sec:further:key_management:policies});	
	\item \textbf{performance monitoring}: examine metrics in logs to assess the performance of the system (e.g., encryption time, latency, utilization of computational resources).
\end{itemize}

The availability of logs also enables internal auditing. In the context of \glspl{kms}, organizations should periodically conduct at least the following 3 types of audit \cite{barker_recommendation_2020_p1}:

\begin{itemize}
	\item \textbf{operational audit}: an organization should ensure that its \gls{kms} is properly configured to operate or continue to operate in compliance with the defined key management policy (see \Cref{sec:further:key_management:policies}). This audit may include:
	\begin{itemize}
		\item verifying the efficacy of the adopted security plan and the corresponding procedures put in place to respond to security incidents;
		\item confirm that suitable \gls{ac} policies are properly defined and enforced;
		\item ensure that users are properly trained and educated;
		\item check the functionality of the components of the \gls{kms} (e.g., see inventory management for keys and metadata in \Cref{sec:key_management:inventory});
	\end{itemize}
	\item \textbf{security audit}: an organization should check that the security mechanisms employed (e.g., see authentication and authorization in \Cref{sec:further:key_management:authn}) are effective and provide the expected level of security --- especially in light of new technological developments and attacks;
	\item \textbf{monitoring audit}: an organization should review the activities concerning the \gls{kms} to detect anomalies, attacks, and activities not compliant with the established operational policies and procedures (see \Cref{sec:further:key_management:policies}).
\end{itemize}



\subsubsection{The Key Management Interoperability Protocol}\label{sec:cloud:kmip}

The \gls{kmip}\footnote{OASIS Key Management Interoperability Protocol (KMIP) TC (\url{https://www.oasis-open.org/committees/tc\_home.php?wg\_abbrev=kmip})}$^,$\footnote{OASIS Key Management Interoperability Protocol (KMIP) TC Wiki (\url{https://wiki.oasis-open.org/kmip})} is an \gls{oasis} standard describing an extensible key management protocol that defines message formats for the manipulation of cryptographic keys on \glspl{kms}. The protocol considers the lifecycle proposed by \gls{nist} (see \Cref{sec:key_management}) for both symmetric and asymmetric keys in a variety of formats. Furthermore, the protocol considers operations for handling digital certificates, keys wrapping, provisioning schemes, and cryptographic computations (e.g., encryption) as well as the management metadata associated with the keys. There exist several commercial \glspl{kms} supporting the \gls{kmip} standard, as well as many other \glspl{kms} --- e.g., those offered by Cloud providers --- which do not support the \gls{kmip} standard.\footnote{Key management - Wikipedia (\url{https://en.wikipedia.org/wiki/Key\_management\#Key\_Management\_Interoperability\_Protocol\_(KMIP)})} Organizations should carefully consider whether to employ \glspl{kms} following the \gls{kmip} standard for enabling interoperability across \glspl{kms}.



\subsubsection{Crypto-Agility}\label{sec:further:key_management:agility}

Crypto-agility --- also known as cryptographic agility --- is a proactive and flexible approach to cryptographic key management that refers to the ability of a system (e.g., a \gls{kms}) to \textbf{rapidly, effortlessly, and seamlessly} change employed cryptographic algorithms, ciphers, and cryptosystems in an \textbf{(at least semi-)automated} fashion, \textbf{without requiring a significant overhaul} of (the design of) the system itself and \textbf{without major service disruptions} \cite{barker_getting_2021}. More concretely, crypto-agility comprises the ability to:
\begin{itemize}
	\item switch cryptographic ciphers easily and efficiently;
	\item update key sizes (and other related parameters) to preserve a sufficient level of security;
	\item rotate and renew keys in a regularly and, most importantly, efficiently;
	\item quickly adapt to standards and compliant requirements published by relevant organizations (e.g., \gls{nist});
	\item continuously monitor the cryptographic landscape for new attacks and vulnerabilities;
	\item thoroughly test new cryptographic algorithms, ciphers, and cryptosystems to assess their impact on the system and avoid service disruptions.
\end{itemize}

In other words, the adoption of crypto-agility ensures a rapid response of a system to evolving security threats (e.g., discovery of vulnerabilities in ciphers and breakthroughs in cryptanalysis), updates in standards, and new compliance requirements and regulations. Crypto-agility is especially relevant in light of the advent of quantum computing, which will render many of the most currently used cryptographic ciphers insecure (see \Cref{sec:further:key_management:pq}).



\subsubsection{Post-Quantum Cryptography}\label{sec:further:key_management:pq}

While traditional computing uses binary transistors to represent states of data through 0's and 1's, quantum computing\footnote{What is Quantum Computing? | IBM (\url{https://www.ibm.com/topics/quantum-computing})} employs the principles of quantum mechanics --- such as superposition and entanglement --- to perform calculations.\footnote{Superposition refers to a combination of states that are ordinarily independently (or exclusive to each other), while entanglement refers to a counter-intuitive quantum phenomenon, where entangled particles behave together as a system in a way that cannot be explained using traditional logic.} Quantum computing uses qubits for representing data in more than 2 states and exploring multiple solutions to a problem at the same time.\footnote{Quantum computing: Some (not so) gruesome details (\url{https://plus.maths.org/content/really-how-do-quantum-computers-work})} This ability grants quantum computers significantly higher computational power, which results in the possibility of executing quantum algorithms (e.g., the Shor's algorithm \cite{shor_algorithms_1994}) capable of quickly solving mathematical problems such as the integer factorization and the discrete logarithm problems --- which are the security foundation of the \gls{rsa}, Paillier, and \gls{ecc} ciphers (see \Cref{app:cryptography:asymmetric:ciphers}). In other words, with the advent of quantum computing, several currently used cryptographic ciphers would become insecure. In particular, quantum computing poses a threat to asymmetric cryptographic ciphers, while currently used symmetric cryptographic ciphers (e.g., \gls{aes}) would still be secure given long enough keys --- although post-quantum symmetric ciphers may still be developed as a precautionary measure.

Quantum computing is still being researched and developed (the main obstacle to building large-scale quantum computers being ``noise'' \cite{fosel_reinforcement_2018}) but, at the time of writing, the most optimists predict that quantum computers capable of breaking current asymmetric cryptographic ciphers will be available in little more than 20 years (i.e., around 2045).\footnote{Post-Quantum Cryptography | CSRC (\url{https://csrc.nist.gov/projects/post-quantum-cryptography})} In any case, to prevent adversaries from storing ciphertexts now and using quantum computing to decrypt them in the future (i.e., ``harvest now, decrypt later'' or ``store-and-break`` attack), the concept of post-quantum cryptography --- also known as quantum-resistant cryptography --- has already emerged. Essentially, post-quantum cryptography concerns with the study of mathematical problems that ($i$) can be used as the basis for new cryptographic algorithms interoperable with existing communications protocols and networks and ($ii$) cannot be (efficiently) solved by either traditional or quantum computers.

In 2016, \gls{nist} initiated a process to solicit, evaluate, and standardize post-quantum asymmetric cryptographic ciphers.\footnote{NIST Asks Public to Help Future-Proof Electronic Information (\url{https://www.nist.gov/news-events/news/2016/12/nist-asks-public-help-future-proof-electronic-information})} That process led to the announcement of 4 post-quantum cryptographic algorithms in 2022 that can be used for either encryption or digital signatures (see \Cref{app:cryptography:asymmetric});\footnote{NIST Announces First Four Quantum-Resistant Cryptographic Algorithms (\url{https://www.nist.gov/news-events/news/2022/07/nist-announces-first-four-quantum-resistant-cryptographic-algorithms})} more in detail:

\begin{itemize}
	\item \textbf{encryption}: CRYSTALS-Kyber algorithm.\footnote{Kyber (\url{https://pq-crystals.org/kyber/index.shtml})};
	\item \textbf{digital signatures}: CRYSTALS-Dilithium,\footnote{https://pq-crystals.org/dilithium/index.shtml (\url{https://pq-crystals.org/dilithium/index.shtml})}
	FALCON,\footnote{https://falcon-sign.info/ (\url{https://falcon-sign.info/})}
	and SPHINCS+.\footnote{https://sphincs.org/ (\url{https://sphincs.org/})} The first algorithm is generally recommended, while the second allows for smaller signatures. Finally, although generally worse than the others, SPHINCS+ is a valuable backup option as it relies on a different mathematical problem. Indeed, CRYSTALS-Dilithium and FALCON (and also CRYSTALS-Kyber) are based on structured lattices, while SPHINCS+ is based on hash functions, and it is thus important in case vulnerabilities or flaws are found in the first 3 algorithms.
\end{itemize}

In August 2023, \gls{nist} released draft standards for 3 of these algorithms; a draft standard for FALCON is expected to be finalized in 2024. It is worth noting that further algorithms are currently under evaluation by \gls{nist}.\footnote{NIST to Standardize Encryption Algorithms That Can Resist Attack by Quantum Computers (\url{https://www.nist.gov/news-events/news/2023/08/nist-standardize-encryption-algorithms-can-resist-attack-quantum-computers})} 


\paragraph*{Relevance of Quantum Computing to Key Management.}

It is important for organizations to start planning for the replacement of hardware, software, and services that employ asymmetric cryptographic ciphers theoretically vulnerable to quantum computing. Indeed, as said in \Cref{sec:further:key_management:pq}, although not currently available, quantum computing still poses a threat to sensitive data in light of the ``harvest now, decrypt later'' attack. Thus, organizations should already reason on whether quantum computing poses a threat to their keys (and data). More in detail, if an organization:
\begin{itemize}
	\item requires a certain key (and the corresponding data) to be secure for a number of years {\tt k}, and
	\item expects to require {\tt i} years to harden its cryptographic infrastructure to be secure against quantum computing, and
	\item forecasts that a large-scale quantum computer will be available in {\tt q} years,
\end{itemize} 

then quantum computing is relevant to such an organization --- for the given key --- if {\tt q < k + i}. As intuitive, post-quantum cryptography is therefore also linked to the concept of crypto-agility (see \Cref{sec:further:key_management:agility}. Unfortunately, it is likely that post-quantum cryptographic ciphers will not be easily pluggable in currently used \glspl{kms} due to differences in, e.g., performance, reliability, key size (some post-quantum cryptographic ciphers have very large keys), error handling, and key establishment \cite{newhouse_migration_2023}. 

Organizations are thus encouraged to define a strategy (often called ``quantum-readiness roadmap'') for planning and actuating the migration toward post-quantum cryptography in their \glspl{kms} in a proactive way.\footnote{Quantum-readiness: Migration to Post-Quantum Cryptography (\url{https://www.nccoe.nist.gov/sites/default/files/2023-08/quantum-readiness-fact-sheet.pdf})} The first step of a quantum-readiness roadmap is the --- not trivial --- discovery and identification of cryptographic algorithms, ciphers and cryptosystems (used either in hardware, software or firmware) vulnerable to quantum computing. Indeed, having a complete inventory of quantum-vulnerable assets, along with eventual interdependencies (e.g., with external systems and networks), allows organizations to start a risk assessment process and define a prioritization for the migration; for more details on the topic of post-quantum cryptography and crypto-agility, we refer the interested reader to \cite{barker_getting_2021,newhouse_migration_2023}.

\remarkBlock{Post-quantum cryptography and the cloud}{Some cloud providers offering \glspl{kms} and related services already propose solutions for post-quantum cryptography. For instance, \gls{aws} offers a hybrid post-quantum key exchange option for \gls{tls} connections toward the \gls{aws} Key Management Service.\footnote{Using hybrid post-quantum TLS with AWS KMS - AWS Key Management Service (\url{https://docs.aws.amazon.com/kms/latest/developerguide/pqtls.html})} Nonetheless, it is important to highlight that these solutions are not based on any standards (indeed, \gls{nist} will publish the first standardized post-quantum cryptographic ciphers in 2024, as said in \Cref{sec:further:key_management:pq}) hence not currently meant for being used in production. Instead, these solutions are offered mainly for testing and early feedback.\footnote{Post-Quantum Crypto - Amazon Web Services (AWS) (\url{https://aws.amazon.com/security/post-quantum-cryptography/})} However, as said in \Cref{sec:cloud:providers}, Cloud services change frequently and quickly, hence Cloud providers could soon propose new (and standardized) solutions for post-quantum cryptography.}


\subsubsection{Inventory Management for Keys and Metadata}\label{sec:further:key_management:inventory}

Throughout its lifecycle, a key is typically associated with metadata such as key identifier, key type (see \Cref{app:cryptography:keys}), intended purpose and corresponding cipher, owner(s) information, cryptoperiod (see \Cref{sec:key_management:lifecycle:operational}), current key state, key history, location of backups, identifiers of \glspl{kek} applied over the key, integrity information, and other relevant information (e.g., key length, requirements and additional conditions for use). When necessary, these metadata should always be transmitted with the corresponding key in any state \cite{barker_profile_2015}. Moreover, organizations should establish and maintain an inventory of (at least long-term) keys (along with the associated metadata) and digital certificates through a dedicated inventory management system \cite{barker_recommendation_2020_p1} managed under a specific key-inventory policy (see \cref{sec:further:key_management:policies}). Key inventories --- which are usually (at least logically) separated from key backups and archives --- can be used to:

\begin{itemize}
	\item identify which entities need to be notified in the event of revocation or rotation and which digital certificates need to be revoked;
	\item monitor the cryptoperiod of keys (and the expiration of digital certificates) for prompt renewal and rotation;
	\item find configurations of algorithms, ciphers, and cryptosystems that are no longer deemed secure (in this regard, see also the concept of crypto-agility in \Cref{sec:further:key_management:agility});
\end{itemize}

\sbnote{TODO
	a
%implement key management practices into a \gls{kms} \cite{barker_recommendation_2019_p2})
}